\section*{Research}
\begin{itemize}
    \item \textbf{Master Thesis}: University of Groningen (2026) \begin{itemize}
        \item \textbf{Title}: \href{https://milenemdvl.github.io/CV/theses/MasterThesis.pdf}{Unstructured Finite Volume Discretisation Scheme for Turbulent Navier-Stokes}
        \item \textbf{Supervisors}: prof. dr. ir. R.W.C.P Verstappen and dr. ir. R. Luppes 
        \item \textbf{Description}: In this paper, I present a Finite Volume discretisation scheme for two- and three-dimensional unstructured grids obtained with Delaunay Triangulation. This triangulation is chosen such that the vertices only lie on the boundaries of the triangles' circumcircles. I first go over the basics of turbulence modelling, such as Large Eddy Simulation and filtering. With the aim of the divergence and gradient being adjoints of one another, I present the derivations for the discrete operators, and I give an example for a simple two-dimensional grid. I solve the pressure via a Poisson problem where the intermediate linear systems are solved iteratively. I verify the code using the method of manufactured solutions, in which I present velocity and density solutions as linear combinations of sines and cosines and aim to recover these solutions.
    \end{itemize}
    \item \textbf{Internship}: Von Karman Insitute for Fluid Dynamics (2026) \begin{itemize}
        \item \textbf{Title}: \href{https://milenemdvl.github.io/CV/theses/InternshipThesis.pdf}{Towards Efficient Galerkin Projection in
        Lagrangian ROMs for Sloshing}
        \item \textbf{Supervisors}: prof. M. A. Mendez, ir. S. Ahizi,  dr. J. Koellermeier, and dr. ir. R.
        Luppes,
        \item \textbf{Description}: In this paper I derive a reduced order model for
        smoothed particle hydrodynamics, by rewriting the standard neighbour
        sums as dense matrix operations, and using proper orthorgonal
        decomposition and Galerkin projection to obtain a linear time invariant
        model. However, the kernel transformations inherent to particle dynamics
        introduce a nonlinear state dependent term. To this end, I evaluate
        three methods: linearisation, random Fourier features, and a singular
        value decomposition based interpolation. 
    \end{itemize}
    \item \textbf{Bachelor Thesis}: University of Groningen (2023) \begin{itemize}
        \item \textbf{Title}: \href{https://milenemdvl.github.io/CV/theses/BachelorThesis.pdf}{Average Surface Roughness for 2-Dimensional Atmospheric Flow}
        \item \textbf{Supervisors}: dr. ir. R. Luppes and dr. A. E. Sterk 
        \item \textbf{Description}: In this paper I present a study of four different methods of computing an
        average (effective) roughness length (following André and Blondin (1986), Taylor
        (1989), de Vries et. al. (2003), and taking the average of the
        logarithmic law).
    \end{itemize}
\end{itemize}